To interpret that group valuation, consider what would be paid to transfer
the liability in an orderly transaction between market participants at the
measurement date. The obligation remains outstanding, and the transferee must
fulfil it; the thought experiment is not an immediate settlement with the
policyholder. AASB~13 therefore calls for assumptions that market participants
would use, in the accessible principal market or, in its absence, the most
advantageous market. The unit being valued still comes from C20: the insurance
group \citep[paragraphs~2--3, 9, 13--24 and 34]{aasb13comp2021}.

A quoted price is not required in every case. The valuation should maximise
relevant observable inputs, while supported unobservable inputs can supply
the participant assumptions for which market observations are unavailable.
The insurer's own data can be a starting point, but require adjustment when
reasonably available evidence indicates different participant assumptions or
entity-specific features unavailable to other participants. A technique must
also fit the circumstances and available data. These requirements make the
source, date and interpretation of each valuation input consequential
\citep[paragraphs~61--67 and 87--89]{aasb13comp2021}.

There is a further reason to retain separate fair-value and FCF calculations.
AASB~13 includes non-performance risk, including the issuer's own credit
risk, in liability fair value and assumes that risk is unchanged by the
hypothetical transfer. AASB~17 expressly excludes the issuer's own
non-performance risk from fulfilment cash flows for issued insurance contracts.
Even when the two measurements draw on some of the same cash-flow records,
their treatment of this risk differs. Neither a single imported FCF balance
nor the value of the underlying investments supplies both C20 inputs
\citep[paragraphs~42--44]{aasb13comp2021}
\citep[paragraph~31]{aasb17}.\footnote{The comparison uses AASB~108's
31 December 2022 compilation, applying to periods beginning from 1 January
2023, and AASB~13's 30 June 2021 compilation, applying from 1 July 2021.
These are retained historical editions; later amendments and the actual
entity's applicable requirements need their own check.}
